\section{Behaviour}\label{sec:behaviour}

\paragraph{Splicing damage grows with depth.} Splicing the reconstruction into GPT-2 at every non-BOS position raises
next-token cross-entropy by 0.009--0.026 nats at layer 1 and by 0.213--0.358 nats at layer 11
(Figure~\ref{fig:recon}, right), while loss recovered relative to mean ablation stays at or above 0.93. An identity
splice changes nothing, which validates the hook.

\paragraph{SAE error is usually more harmful than random error.} Replacing the SAE error by a norm-matched random error
gives a smaller cross-entropy increase in most runs: the ratio of SAE-induced to random-induced damage is 0.76--2.89
per run (median 1.64), and the seven runs below 1 are all at layers 1--3. The SAE error is therefore not behaviourally
neutral, although this control does not identify which directions are responsible.

\paragraph{Within a layer, reconstruction error predicts behaviour but not truth loss.} Pooling SAEs across layers
mixes the depth trend with differences between SAEs; a pooled correlation of 0.89 between reconstruction quality and
persona preservation, for instance, falls to 0.05 within layer. We therefore regressed each outcome on reconstruction
error within layer, with layer and architecture fixed effects and a within-layer permutation test (Appendix
Table~\ref{tab:m0}). Lower OWT NMSE predicts less splice damage ($\beta=0.165$ per unit log NMSE, Holm-corrected
$p=0.0008$; per-layer Spearman $\rho$ 0.80--0.97 from layer 3), and the association survives without architecture
fixed effects. For the frozen truth gap the formal test passes (Holm $p=0.044$), but architecture explains 99.8\,\% of
the within-layer NMSE variance and the association disappears without architecture fixed effects ($p=0.89$).
Reconstruction error measured on the truth statements themselves is not supported as a predictor of either outcome.
Reconstruction error is therefore a robust within-layer predictor of behavioural damage, but not of truth-information
loss.
